\section{Related Work}
\label{sec:related}

Our work bridges two largely separate research streams---RL training with execution feedback and test-time refinement---that have rarely been combined systematically.

\subsection{RL Training for Code Generation}

Reinforcement learning from execution feedback has emerged as a powerful paradigm for code generation. CodeRL \citep{le2022coderl} treats the code-generating language model as an actor and introduces a critic network to estimate functional correctness, using program execution results as reward signals. This approach achieves state-of-the-art results on APPS and HumanEval by optimizing for test-case pass rates rather than token-level likelihood. PPOCoder \citep{shojaee2023ppocoder} extends this framework with proximal policy optimization, demonstrating improved stability during RL fine-tuning. B-Coder \citep{yu2023bcoder} introduces value-based RL as an alternative to policy-based methods, leveraging off-policy programs to improve sample efficiency.

\textbf{Limitation:} These works evaluate RL-trained models in single-shot mode only. Whether RL training improves the model's ability to iteratively refine code based on execution feedback remains untested.

\subsection{Test-Time Refinement}

A parallel line of work focuses on improving code generation at inference time without additional training. Self-Refine \citep{madaan2023selfrefine} demonstrates that large language models can iteratively refine their outputs using self-generated feedback, achieving $\sim$20\% improvement across diverse tasks including code generation. S* \citep{li2025sstar} presents a hybrid test-time scaling framework that enables 3B models to match or exceed GPT-4o-mini on HumanEval through execution-guided search. Recent work on test-time compute scaling \citep{ma2025testtimescaling} shows that 32B models can achieve 46\% on SWE-bench Verified, outperforming DeepSeek R1 671B through strategic test-time allocation.

\textbf{Limitation:} These methods apply refinement to pretrained or instruction-tuned models. Whether RL training provides a better starting point for refinement is not explored.

\subsection{Our Position}

Unlike prior work that studies RL training or test-time refinement in isolation, we provide the first factorial experimental framework to measure their interaction. We do not claim to improve upon CodeRL or Self-Refine individually; rather, we ask whether combining them produces superadditive gains. Our $I(F;E)$ metric provides a mechanistic probe beyond aggregate accuracy, enabling analysis of \emph{why} interaction effects occur (or fail to occur).
